% Notes of the Montreal seminar, 1970 (cote 129, typed part): archive cote 129, site page 62 (PDF page 63).
% typescript. Draft compared: none.
% AI-transcribed by Claude Opus 5.5 (claude-opus-5-5, adaptive thinking, effort high), Message Batches API, batch mtl129_full (msgbatch_01AGnDKHwy6ATCM7ehUa6263), 2026-11-13; not yet checked by eye.
\scanpage{129/62}
\mspage{\hand{III-7}}

4. \uline{Groupes de Barsotti-Tate}

La terminologie de groupe de Barsotti-Tate est préférée à
celle de ``groupe p-divisible'' employée par Tate (J.T. Tate, p-divisi-%
ble groups, in Proceedings of a Conference on Local Fields, Driebergen
1966 (Springer)), la notion de groupe p-divisible étant beaucoup plus
générale:

4.1. \uline{Définition}. On dit qu'un groupe $G$ ($\mathcal{T}$-faisceau en groupes)
est p-divisible si $p.\mathrm{Id}\colon G \to G$ est un épimorphisme.

Si $G$ est p-divisible, $G(n)$ est un $\Lambda_n$-module plat d'après
2.2 (e), et le groupe de p-torsion $G(\infty)$ est lui aussi p-divisible.
Plus généralement, si on se donne un système inductif
\margin{X} \margin{n+1}
\[ G_1 \hookrightarrow G_2 \hookrightarrow \ldots\ G_n \hookrightarrow G_{n\over{\unsure{-}}{+}1} \hookrightarrow \ldots \]
tel que $G_n$ est annulé par $p^n$ et $G_n(n-1) = G_{n-1}$ (on appelle un tel
système de groupes un \uline{système co-p-adique}) et tel que de plus $G_n$ soit
un $\Lambda_n$-module plat, alors $\lim\limits_{\rightarrow} G_n$ est un groupe p-divisible et de p-torsion.

4.2. \uline{Définition}. Un \uline{p-groupe de Barsotti-Tate sur $S$} (ou groupe
de Barsotti-Tate) est un groupe $G$ ($\mathcal{T}$-faisceau en groupes commutatifs)
p-divisible et de p-torsion tel que $G(1)$ soit un schéma fini localement
libre sur $S$.

Il est équivalent de dire que l'un quelconque des $G(n)$ est
un schéma fini localement libre sur $S$ (2.4).

\mspage{\hand{62}}
